\chapter{MRSA Screening}

\section*{What Is This Test?}
MRSA screening detects colonization or infection with methicillin-resistant \textit{Staphylococcus aureus} (MRSA). MRSA is any strain of \textit{S. aureus} that has acquired the \textit{mecA} gene (or the \textit{mecC} homologue), which encodes an altered penicillin-binding protein (PBP2a) with low affinity for all beta-lactam antibiotics (penicillins, cephalosporins, carbapenems), rendering them resistant. MRSA is broadly divided into healthcare-associated MRSA (HA-MRSA), associated with invasive infections including bacteremia, surgical site infections, and ventilator-associated pneumonia, typically in patients with healthcare exposures, indwelling devices, or recent hospitalization; and community-associated MRSA (CA-MRSA), most commonly causing skin and soft tissue infections (abscesses, furunculosis, cellulitis) in otherwise healthy individuals without traditional healthcare risk factors. CA-MRSA strains (typically USA300 clone, pulsed-field type) often carry the Panton-Valentine leukocidin (PVL) virulence factor and are susceptible to more non-beta-lactam antibiotics than HA-MRSA strains. MRSA colonization of the anterior nares is the major risk factor for subsequent MRSA infection. Active detection and isolation (ADI) programs routinely screen high-risk patients for MRSA nasal carriage using culture or PCR. MRSA infections require treatment with non-beta-lactam agents: vancomycin (first-line for serious infections), daptomycin, linezolid, ceftaroline, or clindamycin/trimethoprim-sulfamethoxazole/doxycycline for mild-to-moderate infections.

\section*{Alternative Names}
Methicillin-Resistant \textit{S. aureus} Screening; MRSA Culture; MRSA PCR; MRSA Nasal Swab; MRSA Surveillance; MRSA Screen; MRSA Active Detection Testing

\section*{Principle}
MRSA screening culture: A nasal swab (also axillae, groin, wounds if present) is inoculated onto a selective chromogenic agar (CHROMagar MRSA, BBL CHROMagar Staph aureus, Brilliance MRSA agar, or Spectra MRSA). These media contain cefoxitin (a potent inducer of \textit{mecA} expression) and chromogenic substrates that produce characteristic colony color when cleaved by \textit{S. aureus} enzymes (mauve/pink colonies on CHROMagar MRSA; blue on Spectra MRSA). Plates are incubated at 35--37\textdegree{}C for 24--48 hours. Colonies are confirmed as \textit{S. aureus} by coagulase test (slide or tube), latex agglutination for protein A and clumping factor, or MALDI-TOF. Oxacillin or cefoxitin disk diffusion or broth microdilution confirms resistance with cefoxitin MIC $\geq$4 mcg/mL or oxacillin MIC $\geq$4 mcg/mL for MRSA. MRSA PCR: Real-time PCR directly from nasal swabs detecting the \textit{mecA} gene and an \textit{S. aureus}-specific target (e.g., the \textit{nuc} gene encoding thermonuclease, or the \textit{spa} gene encoding protein A). Platforms: Cepheid GeneXpert MRSA (results in 1 hour), Roche cobas MRSA/SA (results in 2 hours), BD MAX MRSA (results in 2 hours). Sensitivity 93--98\%, specificity 95--99\%. PCR provides results in 1--2 hours compared to 24--48 hours for culture. However, PCR detects \textit{mecA} DNA and cannot distinguish viable MRSA from nonviable organisms or from the co-colonization of methicillin-resistant coagulase-negative staphylococci (CoNS) carrying \textit{mecA} alongside MSSA (a false-positive MRSA result by some older PCR assays; newer assays include \textit{S. aureus}-specific targets to improve specificity). CHROMagar culture remains the gold standard. Susceptibility testing: Confirmed \textit{S. aureus} isolates are tested by broth microdilution, disk diffusion (cefoxitin disk), or automated systems (VITEK, MicroScan, BD Phoenix). Vancomycin MIC (by broth microdilution or gradient diffusion/Etest) should be determined for all MRSA isolates from serious infections given concerns about vancomycin-intermediate \textit{S. aureus} (VISA, MIC 4--8 mcg/mL) and vancomycin-resistant \textit{S. aureus} (VRSA, MIC $\geq$16 mcg/mL, \textit{vanA}-mediated, extremely rare).

\section*{History}
\textit{S. aureus} was first described by Alexander Ogston in 1880. Penicillin resistance emerged rapidly after penicillin introduction in the 1940s (mediated by beta-lactamase). Methicillin (a beta-lactamase-resistant penicillin) was introduced in 1959; MRSA was reported within one year in 1960 (in the UK). The \textit{mecA} gene was sequenced in the 1980s. MRSA became a major nosocomial pathogen in the 1980s--1990s, and CA-MRSA emerged dramatically in the late 1990s (USA300 clone). PCR-based MRSA screening from nasal swabs was introduced in the mid-2000s (Cepheid GeneXpert MRSA FDA-cleared in 2007). The effectiveness of active detection and isolation (ADI) in reducing MRSA transmission has been demonstrated in numerous studies, though universal decolonization (chlorhexidine bathing + intranasal mupirocin for all ICU patients) was shown to be superior to targeted MRSA screening and isolation in the REDUCE MRSA trial (Huang et al., NEJM 2013). Vancomycin remained the gold-standard therapy for serious MRSA infections since the 1980s. Daptomycin (2003), linezolid (2000), ceftaroline (2010), dalbavancin (2014), oritavancin (2014), and telavancin (2009) expanded the therapeutic armamentarium.

\section*{Why This Test Is Ordered}
\begin{itemize}
  \item Active surveillance for MRSA nasal colonization in patients admitted to ICUs, burn units, and transplant units, or patients transferred from long-term acute care facilities (LTACHs) or nursing homes, per institutional infection control policy
  \item Preoperative MRSA screening before high-risk surgical procedures: cardiac surgery (sternotomy, valve replacement, CABG), orthopedic implant surgery (total hip/knee arthroplasty, spine fusion), and neurosurgery with implant placement. MRSA-colonized patients receive a bundle of preoperative decolonization (intranasal mupirocin BID x 5 days + chlorhexidine bathing daily x 5 days) and intraoperative vancomycin prophylaxis in addition to the standard cephalosporin
  \item Diagnosis of MRSA infection in clinical specimens: Wound drainage, abscess aspirate, blood, sputum (ventilator-associated pneumonia), urine, or tissue. \textit{S. aureus} isolated in culture must be tested for methicillin susceptibility
  \item MRSA screening in patients with recurrent \textit{S. aureus} skin and soft tissue infections to identify persistent MRSA colonization and institute a decolonization strategy (mupirocin, chlorhexidine, dilute bleach baths) after acute infection resolves
  \item Outbreak investigation when a cluster of MRSA infections is identified in a hospital unit, to screen patients and potentially healthcare workers for MRSA carriage and identify environmental reservoirs
  \item Screening of dialysis patients for MRSA colonization, as dialysis patients have high rates of MRSA carriage and are at significantly elevated risk of MRSA bloodstream infection (vascular access-related)
  \item MRSA nasal screen for antibiotic stewardship: A negative MRSA nasal screen by PCR has a negative predictive value of >98\% for MRSA pneumonia. In a patient with suspected hospital-acquired pneumonia and a negative MRSA nasal PCR, empiric MRSA coverage (vancomycin, linezolid) can be safely discontinued or withheld, reducing unnecessary vancomycin exposure and nephrotoxicity
\end{itemize}

\section*{Is This a Primary or Secondary Test}
MRSA screening by culture or PCR is a primary surveillance test ordered based on institutional infection control policy or preoperative protocol, not by clinical symptoms. MRSA identification and susceptibility testing from clinical specimens (wound, blood, sputum) is a primary test in patients with suspected infection. The MRSA nasal PCR as an antibiotic stewardship tool (to discontinue empiric vancomycin for pneumonia) is a secondary decision-support test.

\section*{How This Test Is Performed}
MRSA nasal screening: A flocked swab moistened with sterile saline is inserted 1--2 cm into one anterior nare, rotated against the nasal mucosa 5 times, then the same swab is used to sample the other nare. The swab is placed in the appropriate transport medium: Amies transport medium for culture; manufacturer-specific buffer for PCR (dry swab for GeneXpert). For additional sites (axillae, groin, wounds, throat), separate swabs may be pooled (depending on laboratory protocol). Blood cultures: 2 sets (aerobic + anaerobic) from separate venipuncture sites. Wound/abscess culture: Aspirated pus or tissue is preferred over surface swabs, which may be contaminated with colonizing flora. Sputum culture for suspected MRSA pneumonia: Expectorated sputum, endotracheal aspirate, or BAL.

\section*{What Preparation Is Needed}
No special preparation for nasal screening. For patients already receiving intranasal mupirocin or chlorhexidine bathing (decolonization), screening cultures and PCR should be collected before these agents are initiated or at least 48 hours after discontinuation. Recent antibiotic therapy (particularly intravenous vancomycin, linezolid, or daptomycin) may convert screening cultures from positive to negative.

\section*{Contraindications and Precautions}
A negative MRSA nasal screen (PCR or culture) from a single site (anterior nares) has a sensitivity of only 80--90\% for detecting MRSA colonization because 10--20\% of MRSA carriers are colonized at extra-nasal sites (throat, axillae, perineum, groin, rectum, wounds) but not the nares. For maximum sensitivity in high-risk settings, combined nasal, throat, axilla, and wound/device exit-site screening is recommended. PCR detects dead (nonviable) MRSA, so PCR may remain positive for days after successful decolonization. Culture should be used to document clearance. MRSA PCR for pneumonia de-escalation is based on the high NPV (>98\%) of a nasal screen: MRSA pneumonia is extremely unlikely if the nasal screen is negative. However, this strategy has been validated primarily in community-onset pneumonia and in patients NOT already receiving anti-MRSA antibiotics. A patient already on vancomycin for 48--72 hours may develop a negative nasal screen despite having MRSA pneumonia (the nasal screen is suppressed by systemic antibiotics). A positive nasal screen does NOT confirm MRSA pneumonia; the NPV is high but the PPV is low (~30\%), because many patients colonized with MRSA do not have MRSA pneumonia. \textit{S. aureus} isolates with borderline oxacillin MICs (2--4 mcg/mL) due to beta-lactamase hyperproduction (BORSA/borderline oxacillin-resistant \textit{S. aureus}) lack \textit{mecA} and are not MRSA; they may be misinterpreted as MRSA if cefoxitin disk testing alone is used without \textit{mecA} PCR or PBP2a latex agglutination.

\section*{Risks}
Nasal swabbing is painless and risk-free. Risks of blood culture, wound culture, and sputum induction are related to the collection procedure, not the test. The primary risk of MRSA is failure to identify it in a timely manner, leading to inappropriate empiric antibiotic therapy and worse outcomes (MRSA bacteremia mortality: 15--30\% with appropriate therapy; >50\% with delayed/inappropriate therapy).

\section*{What Happens After the Test}
MRSA PCR results in 1--2 hours. Culture results: presumptive positive at 18--24 hours; final negative at 48 hours. MRSA-colonized patients should be placed in contact precautions (gown and gloves for all patient contact) per CDC guidelines. Decolonization with intranasal mupirocin BID x 5 days and chlorhexidine bathing daily x 5 days may be instituted to reduce the risk of subsequent MRSA infection. For MRSA identified from clinical specimens (blood, wound, sputum), antimicrobial susceptibility testing guides therapy. Vancomycin is the empiric antibiotic of choice for serious MRSA infections (target trough 15--20 mcg/mL for invasive infections, with AUC-guided dosing preferred over trough-only monitoring to achieve AUC/MIC $\geq$400--600). Alternative agents are selected based on infection site, severity, MIC results, renal function, and drug tolerability.

\section*{Understanding Results}

\subsection*{Below Normal}
\begin{itemize}
  \item \textbf{MRSA culture negative (no MRSA isolated):} No MRSA detected from the sampled sites. Low risk of MRSA colonization or infection. For nasal screening, the negative predictive value for MRSA pneumonia exceeds 98\%, making it an excellent rule-out test for considering discontinuation of empiric MRSA antibiotic coverage. Negative predictive value for subsequent MRSA clinical infection in hospitalized patients: >95\%. A negative culture does not exclude MRSA colonization at unsampled sites
  \item \textbf{MRSA PCR negative (not detected):} Same interpretation as culture negative, with the added advantage of rapid turnaround (1--2 hours vs. 24--48 hours)
  \item \textbf{Blood culture growing \textit{S. aureus} with oxacillin MIC $\leq$ 2 mcg/mL (MSSA):} Methicillin-susceptible \textit{S. aureus} (MSSA). Preferred therapy: nafcillin/oxacillin 2 g IV q4h or cefazolin 2 g IV q8h. Vancomycin is inferior to beta-lactams for serious MSSA infections and should NOT be used if a beta-lactam can be tolerated
\end{itemize}

\subsection*{Above Normal}
\begin{itemize}
  \item \textbf{MRSA culture positive:} MRSA colonization or infection confirmed. \textit{S. aureus} with oxacillin MIC $\geq$4 mcg/mL and/or cefoxitin disk diffusion zone <22 mm (or cefoxitin MIC $\geq$4 mcg/mL). The isolate should be tested for susceptibility to vancomycin, daptomycin, linezolid, clindamycin (inducible resistance by D-test), trimethoprim-sulfamethoxazole, tetracyclines (doxycycline, minocycline), and in some cases, ceftaroline, rifampin (NEVER as monotherapy; resistance develops rapidly)
  \item \textbf{MRSA blood culture positive:} Diagnostic of MRSA bacteremia. All patients with MRSA bacteremia require: (1) Source identification and control: remove/remove-replace intravascular catheters, drain abscesses, debride infected tissue. (2) Echocardiography: TEE is preferred (sensitivity >95\% for IE) over TTE (sensitivity ~50\%); all patients should be evaluated for infective endocarditis (IE). (3) Infectious diseases consultation is recommended and associated with significantly reduced mortality. (4) Duration of therapy: uncomplicated bacteremia (sterile blood cultures within 72 hours, no endocarditis, no implants, defervescence within 72 hours) = 14 days IV; complicated bacteremia (endocarditis, persistent bacteremia, metastatic foci) = 4--6 weeks IV
  \item \textbf{Vancomycin MIC elevated (2 mcg/mL):} "High" vancomycin MIC within the susceptible range. Associated with higher rates of treatment failure when vancomycin is used, particularly with trough-only dosing. AUC-guided dosing is now recommended over trough-only monitoring. Alternative agents (daptomycin, linezolid) should be considered for serious infections with vancomycin MIC $\geq$2 mcg/mL
  \item \textbf{D-test positive (inducible clindamycin resistance):} Erythromycin-resistant, clindamycin-susceptible on routine testing, but a positive D-test (flattening of the clindamycin zone adjacent to the erythromycin disk) indicates inducible \textit{erm} gene methylation conferring resistance to macrolides, lincosamides, and streptogramin B (MLS$_{B}$ phenotype). Clindamycin should NOT be used for serious infections despite apparent susceptibility, as constitutive resistance can emerge during therapy with clinical failure
\end{itemize}

\section*{Normal Results}
\textbf{No MRSA isolated} or \textbf{MRSA not detected} from nasal screening. Blood/wound/sputum cultures: \textbf{No \textit{Staphylococcus aureus} isolated} or \textbf{MSSA isolated} (methicillin-susceptible).

\section*{Follow-up Tests}
\begin{itemize}
  \item \textbf{Vancomycin MIC determination (broth microdilution or Etest):} For all MRSA isolates from serious infections; elevated MICs (2 mcg/mL) predict reduced vancomycin efficacy and warrant consideration of alternative agents
  \item \textbf{Daptomycin susceptibility:} For MRSA from serious infections when daptomycin is being considered; daptomycin is inactivated by pulmonary surfactant and cannot be used for pneumonia
  \item \textbf{Blood cultures (repeat):} Daily blood cultures should be obtained for patients with MRSA bacteremia to document clearance (sterile within 48--96 hours expected); persistent bacteremia beyond 72--96 hours suggests an uncontrolled source, endocarditis, or other metastatic focus
  \item \textbf{Echocardiography (TEE preferred):} All patients with MRSA bacteremia require evaluation for endocarditis; community-acquired or healthcare-associated with risk factors for IE (prosthetic valve, pacemaker, prolonged bacteremia, embolic phenomena, injection drug use) should have TEE
  \item \textbf{Decolonization confirmation culture:} If decolonization is performed (mupirocin + chlorhexidine) for patients with recurrent MRSA infections, nasal (and possibly extra-nasal) cultures 3--7 days after completion of the decolonization regimen can document clearance. However, recolonization with the same or a new MRSA strain occurs in 30--50\% within weeks to months
\end{itemize}

\section*{References}
\begin{enumerate}
  \item Liu C, Bayer A, Cosgrove SE, et al. Clinical practice guidelines by the IDSA for the treatment of methicillin-resistant \textit{Staphylococcus aureus} infections in adults and children. Clin Infect Dis. 2011;52(3):e18-e55.
  \item Huang SS, Septimus E, Kleinman K, et al. Targeted versus universal decolonization to prevent ICU infection. N Engl J Med. 2013;368(24):2255-2265.
  \item Calfee DP, Salgado CD, Milstone AM, et al. Strategies to prevent methicillin-resistant \textit{Staphylococcus aureus} transmission and infection in acute care hospitals: 2014 update. Infect Control Hosp Epidemiol. 2014;35(7):772-796.
  \item Holland TL, Arnold C, Fowler VG. Clinical management of \textit{Staphylococcus aureus} bacteremia: a review. JAMA. 2014;312(13):1330-1341.
  \item Parente DM, Cunha CB, Mylonakis E, Timbrook TT. The clinical utility of methicillin-resistant \textit{Staphylococcus aureus} (MRSA) nasal screening to rule out MRSA pneumonia. Clin Infect Dis. 2018;67(1):1-7.
\end{enumerate}

\clearpage
